\section{Evaluation}
\label{sec:evaluation}

\sloppy

We implement \Name with $\sim$10K LoC in Python. 
The FusedLoRA and FusedMultiLoRA kernels are developed using Triton~\cite{tillet2019triton}, and multi-adapter pipeline parallelism is built on top of Megatron-LM~\cite{megatron-1}.
Since Megatron-LM does not natively support LoRA, we integrate Hugging Face Transformers~\cite{hf-transformers} and the PEFT Library~\cite{peft} for model architecture and LoRA adaptation.

The optimizations in \Name are designed to be lossless, guaranteeing they do not affect model convergence or final quality.
Our FusedLoRA and FusedMultiLoRA kernels are numerically stable, producing outputs that are functionally identical to the baseline implementations within numerical precision.
While our adaptive scheduler rearranges samples to form balanced microbatches, it strictly preserves the order of global batches, ensuring the sequence of gradient updates remains unchanged. 
Given that model behavior is identical by design, our evaluation focuses exclusively on system performance metrics like throughput.

We evaluate \Name across a range of real-world fine-tuning scenarios involving multiple datasets, model scales, GPU platforms, and job configurations.
Our primary metric is throughput, measured in trained tokens per second, which better reflects system efficiency for inputs with sequence length variations.
Finally, we perform detailed scalability studies, ablation studies, and performance breakdowns to analyze the contribution of each system component.

\begin{figure}[t]
    \centering
    \includegraphics[width=0.975\linewidth]{figures/evaluation-fig-0-data-statistics.pdf}
    \caption{Distribution of sample lengths across the XSum~\cite{xsum-dataset}, CNN/DailyMail~\cite{cnn-dailymail-dataset}, and WikiSum~\cite{wikisum-dataset} datasets used for LoRA fine-tuning.}
    \label{fig:evaluation-fig-0-data-statistics}
\end{figure}


\subsection{Methodology}
\label{sec:evaluation-methodology}

\myparagraph{Hardware Settings} 
We primarily benchmark \Name on a GPU cluster with NVIDIA H100 (80GB) GPUs and additionally report results on L40S (48GB) GPUs to demonstrate generalizability.
Each H100 node is equipped with 8$\times$H100 GPUs connected via NVLink, 208 vCPUs, and Infiniband for multi-node communication. 
Each L40S server contains 4$\times$L40S GPUs connected over PCIe and 128 vCPUs. 
Most experiments use the smallest number of GPUs that fit the model and maintain good utilization, typically 1, 2, or 4 GPUs. 
As shown in our Scalability Studies (Section~\ref{sec:evaluation-scalability-studies}), assigning fewer GPUs per job and using additional GPUs to run more independent jobs often leads to better efficiency, as it reduces inter-GPU communication and synchronization overhead.
The Scalability Studies also show \Name is fully compatible with both data-parallel and multi-node scaling.

\myparagraph{Workload Settings} 
We evaluate \Name on three open-source language models of varying sizes: LLaMa-3.1-8B~\cite{llama3.1}, Qwen-2.5-32B~\cite{yang2024qwen2-5}, and LLaMa-3.1-70B~\cite{llama3.1}. 
All experiments use summarization as a representative sequence-to-sequence task, which is widely used in prior LoRA fine-tuning studies~\cite{hu2022lora, li2023loftq, liu2024rst}.
We select three public summarization datasets: XSum~\cite{xsum-dataset}, CNN/DailyMail~\cite{cnn-dailymail-dataset}, and WikiSum~\cite{wikisum-dataset}. These datasets have diverse length distributions, as shown in Figure~\ref{fig:evaluation-fig-0-data-statistics}, which stresses batching and scheduling under realistic conditions.
For multi-LoRA experiments, we train four LoRA adapters in parallel. 
In the XSum, CNN/DailyMail (CNNDM), and WikiSum configurations, all four adapters are trained independently on the same dataset. In the Mixed setting, each adapter is trained on a dataset combining samples from all three. In the Heterogeneous (Het) setting, the four adapters are trained on different datasets: one each on XSum, CNN/DailyMail, WikiSum, and Mixed.

\myparagraph{Baselines} 
We compare \Name against three baselines: (i) Megatron-LM~\cite{megatron-1} with fully sharded data parallel (FSDP), (ii) Megatron-LM with pipeline parallelism (PP), and (iii) mLoRA~\cite{ye2023mlora}. 
Megatron-LM does not support multi-LoRA fine-tuning natively, so tasks are trained sequentially, while mLoRA supports multi-LoRA fine-tuning.
The original mLoRA uses Python RPC for inter-GPU communication, which performs poorly on NVLink-equipped GPUs. 
Therefore, we reimplement mLoRA inside our system with high-performance communication primitives to ensure fair comparison.
In addition, since mLoRA does not provide a unique multi-LoRA CUDA kernel, we optimistically assume it has the same performance as the naive single LoRA kernel.
Tensor parallelism is not evaluated due to the lack of efficient support in existing LoRA frameworks. 
All experiments use PyTorch 2.6, CUDA Toolkit 12.4, Triton 3.2.0, and Megatron-Core 0.11.0.

\subsection{End-to-End Results}
\label{sec:evaluation-end-to-end-results}

\begin{figure*}[t]
    \centering
    \includegraphics[width=0.99\textwidth]{figures/evaluation-fig-1-end-to-end.pdf}
    \caption{End-to-end training throughput (tokens/sec) of training 4 LoRA adapters on 1, 2, and 4 H100 GPUs. The first four bars per subfigure represent homogeneous workloads (same dataset), and the final (Het) shows heterogeneous adapters trained on different datasets.}
    \label{fig:evaluation-fig-1-end-to-end}
\end{figure*}

\begin{figure}[t]
    \centering
    \includegraphics[width=0.975\linewidth]{figures/evaluation-fig-2-l40-gpus.pdf}
    \caption{End-to-end training throughput (tokens/sec) of training 4 LoRA adapters on 1 and 4 L40S GPUs.}
    \label{fig:evaluation-fig-2-l40-gpus}
\end{figure}

\myparagraph{Speedup on H100 GPUs} 
Figure~\ref{fig:evaluation-fig-1-end-to-end} reports the end-to-end throughput of training 4 LoRA adapters across three models. \Name consistently outperforms all baselines by $1.19-1.96\times$.
% LLaMa-3.1-8B
For LLaMa-3.1-8B, which fits on a single H100 GPU, \Name achieves an average 1.26$\times$ speedup (up to 1.43$\times$), primarily from the FusedLoRA kernel, which reduces memory traffic. 
Since single-GPU setups do not suffer from load imbalance, the improvement here directly reflects kernel-level gains. \Name achieves high speedup on the WikiSum dataset due to the large variance in sample lengths. While the baseline methods suffer from out-of-memory errors, \Name achieves stable packing.
% Other
For Qwen-2.5-32B and LLaMa-3.1-70B, which require distributed training, \Name achieves 1.42$\times$ and 1.64$\times$ average speedup (up to 1.64$\times$ and 1.81$\times$) respectively. Larger models benefit more from improved scheduling, as pipeline stalls and load imbalance become more pronounced at higher parallelism. In the most challenging heterogeneous setting (Het), where each adapter uses a different dataset, \Name still achieves strong performance, highlighting its robustness.

\begin{sloppypar}
\myparagraph{Speedup on L40S GPUs} 
Figure~\ref{fig:evaluation-fig-2-l40-gpus} presents results on NVIDIA L40S GPUs. \Name achieves  $1.19-1.91\times$ average speedup for LLaMa-3.1-8B and Qwen-2.5-32B respectively. The benefit is smaller for LLaMa-3.1-8B due to limited memory capacity on a single L40S GPU, which constrains batch size and limits kernel fusion effectiveness. However, even under such constraints, \Name maintains consistent improvements, demonstrating generalizability across model sizes and hardware platforms.
\end{sloppypar}

\begin{figure}[t]
    \centering
    \includegraphics[width=0.976\linewidth]{figures/evaluation-fig-3-scaling.pdf}
    \caption{Scalability of \Name across 4, 8, and 16 H100 GPUs when training 4 LoRA adapters simultaneously. DP scaling means the more GPUs are used to increase the DP degree for the same job, while Job scaling means different LoRA fine-tuning jobs are scheduled to utilize more GPUs. Global batch sizes are scaled proportionally with GPU count to ensure fair comparison.}
    \label{fig:evaluation-fig-3-scaling}
\end{figure}

\subsection{Scalability Studies}
\label{sec:evaluation-scalability-studies}
We evaluate \Name on 4, 8, and 16 H100 GPUs under two scaling strategies: DP scaling (more GPUs per job) and job-level scaling (more concurrent jobs). 16 H100 GPUs experiment is conducted on 2 nodes each with 8 GPUs, connected via InfiniBand~\cite{infiniband}. Global batch size is scaled with GPU count in both settings.
We draw two key conclusions.
First, job-level scaling consistently outperforms DP scaling due to better load balance, achieving 1.18$\times$ and 1.25$\times$ higher throughput on 8 and 16 GPUs, respectively. 
Second, \Name is fully compatible with DP scaling and multi-node fine-tuning, and still delivers strong performance, achieving 1.78$\times$ average speedup over Megatron-LM and 1.50$\times$ over mLoRA under DP scaling.


\subsection{Effectiveness of FusedLoRA Kernel}
\label{sec:evaluation-effectiveness-of-fused-lora-kernel}

\begin{figure}[t]
    \centering
    \includegraphics[width=1.00\linewidth]{figures/evaluation-fig-5-kernel-perf-fwd_bwd.pdf}
    \caption{Performance of FusedLoRA kernel in forward and backward passes.}
    \label{fig:evaluation-fig-5-kernel-perf-fwd_bwd}
\end{figure}

\begin{figure}[t]
    \centering
    \includegraphics[width=1.00\linewidth]{figures/evaluation-fig-6-layer-perf-normalized.pdf}
    \caption{Performance of FusedLoRA kernel in decoder layers of different models.}
    \label{fig:evaluation-fig-6-layer-perf-normalized}
\end{figure}

\myparagraph{Kernel Performance} 
Figure~\ref{fig:evaluation-fig-5-kernel-perf-fwd_bwd} shows the throughput of our FusedLoRA and FusedMultiLoRA kernels compared to the standard Torch LoRA implementation~\cite{peft}. FusedLoRA achieves an average speedup of $1.27\times$ (up to $1.39\times$), while FusedMultiLoRA achieves $1.17\times$ on average (up to $1.24\times$). In the forward pass, FusedMultiLoRA performs similarly to FusedLoRA, as most computation is shared. In the backward pass, it incurs slight overhead from accumulating gradients across adapters and additional element-wise operations. Despite this overhead, both kernels consistently outperform the baseline across different token sizes and model configurations.

\myparagraph{Layer-wise Performance} 
Figure~\ref{fig:evaluation-fig-6-layer-perf-normalized} compares the speedup across different linear layers in the model. 
FusedLoRA achieves an average speedup of $1.21\times$ (up to $1.30\times$), while FusedMultiLoRA achieves $1.13\times$ (up to $1.17\times$). 
These results are based on microbatches containing four adapters. In practical fine-tuning workloads, each microbatch typically contains only one or two adapters, making FusedMultiLoRA's performance close to FusedLoRA.

\myparagraph{Memory Traffic Reduction}
Figure~\ref{fig:evaluation-fig-7-kernel-ncu-profile} shows DRAM read and write traffic from NVIDIA Nsight Compute (NCU) across representative GEMM shapes. Both FusedLoRA and FusedMultiLoRA consistently reduce memory usage compared to Torch LoRA. For example, on the $8192 \times 4096 \times 4096$ shape, total DRAM traffic reduces to $0.63\times$. Across all settings, traffic is reduced by $34\%-37\%$, confirming that our fusion design effectively reduces redundant memory access.

\myparagraph{Performance Insights Across Diverse Hardware}
The FusedLoRA and FusedMultiLoRA kernels reduce redundant memory access for large activation tensors, which is especially important on hardware where memory bandwidth is much lower compared to compute FLOPS. 
As modern accelerators increase compute FLOPS faster than memory bandwidth~\cite{gholami2024ai-and-memory-wall}, the benefits of our fused kernels are expected to grow in future systems.

\begin{figure}[t]
    \centering
    \includegraphics[width=0.99\linewidth]{figures/evaluation-fig-7-kernel-ncu-profile.pdf}
    \caption{GPU DRAM memory traffic comparison between different kernels from NVIDIA Nsight Compute (NCU).}
    \label{fig:evaluation-fig-7-kernel-ncu-profile}
\end{figure}

\subsection{Effectiveness of Job-Level Scheduling}
\label{sec:evaluation-effectiveness-of-job-level-scheduling}

\begin{figure}[t]
    \centering
    \includegraphics[width=0.99\linewidth]{figures/evaluation-fig-8-pipeline-bubbles.pdf}
    \caption{Pipeline bubble ratio under different methods.}
    \label{fig:evaluation-fig-8-pipeline-bubbles}
\end{figure}

\myparagraph{Pipeline Bubble Reduction} 
Figure~\ref{fig:evaluation-fig-8-pipeline-bubbles} shows how \Name helps reduce pipeline bubbles by scheduling multiple adapters together.
We make three key observations.
First, with only one adapter, the bubble ratio remains high at 44.17\%, close to Megatron-LM's 48.79\%. 
This is because grouping is ineffective when only one dataset is available, showing the importance of multi-LoRA for improved scheduling flexibility.
Second, as more adapters are trained together, the bubble ratio steadily decreases: 15.00\% for 2 adapters, 12.23\% for 3, and 11.09\% for 4. In comparison, mLoRA reaches only 34.11\%, confirming that \Name's grouping and batching significantly reduce pipeline idle time.
Lastly, with four adapters, the bubble ratio is 11.09\%. 
This is due to uneven execution times across pipeline stages, with the last stage taking longer because it handles an extra linear layer and cross-entropy loss.
This limitation is not solvable by our scheduler and is out of scope for this work.

\myparagraph{Tuning Time} 
Figure~\ref{fig:evaluation-fig-9-tuning-time} shows how tuning and computation time grow with the number of training samples for a 4-stage pipeline with 4 adapters, measured on 64 vCPUs and 4 H100 GPUs.
The scheduling time increases nearly linearly, from 15.74 seconds at 640 samples to 102.12 seconds at 25600 samples, demonstrating linear scalability of our scheduler.
The computation time also increases nearly linearly, with a much larger slope than the scheduling time.
The scheduling overhead is negligible for three reasons.
First, the CPU-based scheduling runs in parallel with GPU training of the preceding global batch, with linear scaling of ~4ms per sample in CPU and magnitude difference in execution time between CPU and GPU, making the scheduler's latency fully hidden by this overlap.
Second, as shown in Figure~\ref{fig:evaluation-fig-8-pipeline-bubbles}, the performance gains saturate at 4 adapters, allowing practical deployment with a small constant number of adapters.
Third, we implement a timeout on the MILP solver and fall back to a greedy bin-packing algorithm if the solver takes too long, allowing us to configure the scheduler to balance effectiveness and efficiency, ensuring the scheduling overhead is always within a controllable range.

\begin{figure}[t]
    \centering
    \includegraphics[width=0.99\linewidth]{figures/evaluation-fig-9-tuning-time-with-computation-time.pdf}
    \caption{Tuning and computation time vs. number of samples for 4-stage pipeline with 4 adapters.}
    \label{fig:evaluation-fig-9-tuning-time}
\end{figure}

\myparagraph{Effectiveness of the Merging \& Greedy Fallback}
We evaluate the effectiveness of our scheduler's merging and greedy fallback components on 4 adapters of LLaMa-3.1-70B fine-tuned on four H100 GPUs.
The merging pass improves throughput by 4.34\%, while the two-stage MILP optimization provides an additional 3.82\% improvement over pure greedy bin-packing.
The MILP solver path is selected for 77.4\% of global batches with a timeout of 10 seconds, indicating its effectiveness in reducing token counts for underfilled microbatches.
These modest improvements reflect that most microbatches are already well-packed, and our algorithms primarily optimize the final microbatch in each global batch.
Since scheduling overhead is hidden by parallel GPU execution, these optimizations push performance toward the hardware limit without introducing additional overhead.

\subsection{Speedup Breakdown}
\label{sec:evaluation-speedup-breakdown}

\begin{figure}[t]
    \centering
    \includegraphics[width=0.99\linewidth]{figures/evaluation-fig-4-speedup-breakdown.pdf}
    \caption{Speedup breakdown of \Name on LLaMa-3.1-70B with 4 GPUs.}
    \label{fig:evaluation-fig-4-speedup-breakdown}
\end{figure}

Figure~\ref{fig:evaluation-fig-4-speedup-breakdown} shows the contribution of each component in \Name. 
Starting from the baseline one forward one backward (1F1B) pipeline parallelism used in Megatron-LM, adding FusedLoRA alone yields a 1.13$\times$ speedup. 
This modest gain is constrained by load imbalance and suboptimal token shapes, which limit kernel efficiency (see Figure~\ref{fig:evaluation-fig-5-kernel-perf-fwd_bwd}). 
Replacing 1F1B with Multi-LoRA zero-bubble pipeline parallelism improves throughput to 1.50$\times$ by eliminating pipeline stalls through more microbatches from independent adapters. 
Adding FusedMultiLoRA kernel further raises the speedup to 1.72$\times$ by enabling multi-adapter microbatches and reducing redundant memory access. 
When we apply our scheduler to rebalance token distribution across microbatches, performance improves to 1.57$\times$ even without fusion, as it significantly reduces load imbalance. 
Finally, combining adaptive scheduling with fused kernels achieves the highest speedup of 2.05$\times$, showing the importance of jointly optimizing kernel efficiency, parallelism, and workload balance.

The speedup over mLoRA is driven by two main optimizations: (i) our kernel fusion yields a 1.15$\times$ speedup, as seen by comparing bars 3 and 4 in Figure~\ref{fig:evaluation-fig-4-speedup-breakdown}, with even greater gains when the sequence length is regular and matches the performant sequence length of our kernel (bars 5 and 6); and (ii) our adaptive batching mitigates load imbalance, providing a 1.19$\times$ speedup (bars 4 and 6). These improvements are further supported by microbenchmarks in Figure~\ref{fig:evaluation-fig-5-kernel-perf-fwd_bwd} (1.17$\times$ average speedup, up to 1.24$\times$ for kernel performance) and Figure~\ref{fig:evaluation-fig-8-pipeline-bubbles} (23.02\% reduction in pipeline bubbles).
